% !TeX root = ../../../poisson-cubic.tex
\section{Poisson構造の空間とモジュライスタック}
\label{sec:pg-poisson-stacks}

全ての対象を$\mathbb C$上で考える。Poisson構造のスカラーは割らず、零構造も含める。
まず超曲面を固定し、次に超曲面を動かす場合を扱う。

\subsection{固定した3次超曲面}
\begin{definition}\label{def:pg-pois-space}
滑らかな3次超曲面$X$のPoissonスキームを、2次写像の零スキーム
\begin{equation}\label{eq:pg-pois-space}
 \operatorname{Pois}(X)=Z\left(
 H^0(X,\bigwedge^2T_X)\longrightarrow H^0(X,\bigwedge^3T_X),
 \quad\pi\longmapsto[\pi,\pi]\right)
\end{equation}
と定める。
\end{definition}

\begin{definition}\label{def:pg-fixed-pois-stack}
$\operatorname{Aut}(X)$を$X$の自己同型群スキームとし、
$\phi\cdot\pi=\phi_*\pi$によって$\operatorname{Pois}(X)$に作用させる。
固定モジュライスタックを
\begin{equation}\label{eq:pg-fixed-pois-stack}
 \mathfrak{Pois}(X)=[\operatorname{Pois}(X)/\operatorname{Aut}(X)]
\end{equation}
と定める。点$\pi$の安定化子は
\begin{equation}\label{eq:pg-poisson-aut}
 \operatorname{Aut}(X,\pi)
 =\{\phi\in\operatorname{Aut}(X):\phi_*\pi=\pi\}
\end{equation}
である。
\end{definition}

\begin{proposition}\label{prop:pg-fixed-stack-algebraic}
$n\ge3$の滑らかな3次超曲面$X$について、
$\operatorname{Aut}(X)$は有限エタールであり、
$\mathfrak{Pois}(X)$は有限型Deligne--Mumfordスタックである。また
\begin{equation}\label{eq:pg-fixed-stack-components}
 \operatorname{Irr}(\mathfrak{Pois}(X))
 \simeq\operatorname{Irr}(\operatorname{Pois}(X))/\operatorname{Aut}(X).
\end{equation}
\end{proposition}

\begin{proof}
自己同型群に関する主張はMatsumura--Monsky\cite{PMatsumuraMonsky}の定理の特例である。
ここではgradientのsyzygyから証明する。
Grothendieck--Lefschetzの定理\cite{PGLefschetz}より
$\operatorname{Pic}(X)=\mathbb Z[\mathcal O_X(1)]$である。
従って自己同型は$\mathcal O_X(1)$を保ち、
$H^0(\mathbb P(V),\mathcal O(1))\simeq H^0(X,\mathcal O_X(1))$を通じて
射影線形変換に延長する。
ゆえに$\operatorname{Aut}(X)$は$\mathrm{PGL}(V)$内の$[F]$の閉安定化子である。

そのLie代数を計算する。Euler列を$-3$で捻ると
$H^0(T_{\mathbb P(V)}(-3))=H^1(T_{\mathbb P(V)}(-3))=0$なので、
$H^0(T_{\mathbb P(V)})\simeq H^0(T_{\mathbb P(V)}|_X)$である。
$X$に接する無限小射影変換は、1次式$\ell_i$と定数$c$により
\begin{equation}\label{eq:pg-stack-infinitesimal}
 \sum_i\ell_iF_i=cF
\end{equation}
と表される。Euler等式を引けば
$\sum_i(\ell_i-cz_i/3)F_i=0$となる。
補題\ref{lem:pg-gradient}による1次係数のsyzygyの消滅から
$\ell_i=cz_i/3$であり、このベクトル場は射影化すると零になる。
従って$\operatorname{Lie}\operatorname{Aut}(X)=0$である。
標数0の有限型群スキームは滑らかなので、
$\operatorname{Aut}(X)$は0次元のアフィン群スキーム、従って有限エタールである。

有限型アフィンスキーム$\operatorname{Pois}(X)$の有限エタール群による商は
有限型Deligne--Mumfordスタックである
\cite[Theorem~97.17.2, Tag~06FI]{PStacks}。
商へのアトラスは有限エタール全射であり、
一つの成分の像の逆像はその群軌道に属する成分の和である。
従って相異なる成分の像が一致するのは同じ群軌道に属するときに限り、
\eqref{eq:pg-fixed-stack-components}を得る。
\end{proof}

\subsection{超曲面を動かす相対モジュライ}
$N=n+2$、$V=\mathbb C^N$とし、
$U_N\subset\mathbb P(\operatorname{Sym}^3V^*)$を滑らかな3次式の開集合とする。
普遍超曲面を$p:\mathscr X_N\to U_N$、tautological lineを
$L=\mathcal O_{U_N}(-1)$と書く。
命題\ref{prop:pg-represent}の相対形は
\begin{equation}\label{eq:pg-relative-pois-bundle}
 \mathscr E_N:=p_*\bigwedge^2T_{\mathscr X_N/U_N}
 \simeq\bigwedge^3V\otimes L
\end{equation}
である。$[F]$上では$P\otimes F$を$\iota_{dF}P$に送る。
局所代表を$F$から$aF$へ替えると$P$は$a^{-1}P$へ替わるので、
この対応は重なりで一致する。

この表示とJacobi条件が非被約な基底への変更とも両立することを確認する。
$U_N$のアフィン開集合上で$F$の代表を選ぶと、各次数のgradient Koszul複体は
有限自由加群の複体となる。全ての幾何学的ファイバーで偏微分は2次正則列をなし、
Jacobian環のHilbert級数は
\begin{equation}\label{eq:pg-stack-gradient-hilbert}
 \operatorname{Hilb}_{\mathbb C[z]/(F_1,\ldots,F_N)}(t)=(1+t)^N
\end{equation}
である。従って正のKoszulホモロジーは全ファイバーで零で、
次数0のホモロジーの次元も一定である。
各複体の右端から順に次元を数えると全微分の階数が一定となる。
$U_N$は被約なので、微分を表す行列の所要の小行列式を可逆にした開集合では
像と核が直和因子となる。よって各次数の複体は局所的に分裂し、
任意の底変更後も完全性を保つ。
射影空間のEuler列に用いる線束のコホモロジーも任意の底変更と両立する。
従って命題\ref{prop:pg-represent}の表示、および
命題\ref{prop:pg-jacobi-extension}のEuler・divergence・低次数Koszul完全性による
Jacobi持上げは、任意の$U_N$-スキーム上で成立する。
ここで$F$の各次数の乗法写像も全ファイバーで単射なので局所分裂し、
任意の底変更後に単射である。
従って$F$による約分は可能であり、
式\eqref{eq:pg-euler-jacobi-divisibility}は
$0\to\OO(1)\xrightarrow{F}\OO(4)\to\OO_X(4)\to0$
の大域切断から得られる。
後者で現れる非零整数による除法は$\mathbb C$上で可逆である。

そこで、$\operatorname{Tot}(\mathscr E_N)$上の相対Jacobi零スキームを
\begin{equation}\label{eq:pg-relative-pois-space}
 \mathscr P_n=
 \{([F],\pi):[\pi,\pi]=0\}
 \subset\operatorname{Tot}(\mathscr E_N)
\end{equation}
と定める。局所代表$F$を選べば、これは
$[\iota_{dF}P,\iota_{dF}P]=0$の多項式係数によって定まり、
$P$について2次の有限個の方程式である。
上の底変更の議論により、各ファイバーは非被約構造も含めて
$\operatorname{Pois}(X_F)$と一致する。
代表の変更で括弧は変わらないため、これらの閉スキームは貼り合わさる。

\begin{definition}\label{def:pg-relative-pois-stack}
相対モジュライスタックを
\begin{equation}\label{eq:pg-relative-pois-stack}
 \mathfrak{Pois}^{\mathrm{rel}}_n
 =[\mathscr P_n/\mathrm{PGL}(V)]
\end{equation}
と定める。ここで$\mathrm{GL}(V)$のスカラー行列は
$\bigwedge^3V$と$L$に逆の重みで作用するため、作用は
$\mathrm{PGL}(V)$へ降りる。
このスタックは滑らかな3次超曲面とPoisson構造の組を、
Poisson構造を保つ同型で分類する。
\end{definition}

\begin{theorem}\label{thm:pg-relative-stack-algebraic}
$\mathfrak{Pois}^{\mathrm{rel}}_n$は有限型Deligne--Mumfordスタックであり、忘却射
\begin{equation}\label{eq:pg-relative-forget}
 \mathfrak{Pois}^{\mathrm{rel}}_n\longrightarrow
 \mathfrak{Cub}_n:=[U_N/\mathrm{PGL}(V)]
\end{equation}
は有限型アフィン射である。
$X$を指定した点$\operatorname{Spec}\mathbb C\to\mathfrak{Cub}_n$に沿う
二ファイバーは$\operatorname{Pois}(X)$であり、
その同型類の残余gerbeに沿うファイバーは$\mathfrak{Pois}(X)$である。
\end{theorem}

\begin{proof}
$\mathscr P_n$は有限型スキーム、$\mathrm{PGL}(V)$は滑らかな有限型代数群なので、
商は有限型代数的スタックである
\cite[Theorem~97.17.2, Tag~06FI]{PStacks}。
幾何学的点$(X,\pi)$の安定化子は有限エタール群
$\operatorname{Aut}(X)$の閉部分群であり、標数0なので有限エタールである。
従って対角は不分岐で、商はDeligne--Mumfordである。
忘却射をアトラス$U_N\to\mathfrak{Cub}_n$で底変更すると
$\mathscr P_n\to U_N$を得る。
この射はベクトル束の閉部分スキームの射影であるから有限型アフィンであり、
同じ性質が忘却射へ降下する。
二つのファイバーの記述は、前者では$X$への同型を選び、
後者ではその選択を$\operatorname{Aut}(X)$で割ることによる。
\end{proof}

\begin{remark}\label{rem:pg-stack-component-warning}
連結群$\mathrm{PGL}(V)$は$\mathscr P_n$の各既約成分を保つので、
$\mathfrak{Pois}^{\mathrm{rel}}_n$と$\mathscr P_n$の既約成分は一対一に対応する。
固定した$X$の有限群$\operatorname{Aut}(X)$による商では、
式\eqref{eq:pg-fixed-stack-components}のように成分が同一視され得る。
\end{remark}
